\documentclass[11pt]{article}
\usepackage[margin=1in]{geometry}
\usepackage{booktabs}
\usepackage{amsmath}
\usepackage{graphicx}
\usepackage{xcolor}
\usepackage{hyperref}
\hypersetup{colorlinks=true, linkcolor=blue, urlcolor=blue, citecolor=blue}
\usepackage{enumitem}

\title{\textbf{Edge LLM Inference Across Commodity Accelerators:\\
Measured AMD, NVIDIA, and Apple Performance, and\\
hanzo-engine at llama.cpp Decode Parity on Every Backend}}
\author{Hanzo AI --- hanzo-engine technical white paper, 05}
\date{Measurement pass: 2026-06-17}

\begin{document}
\maketitle

\begin{abstract}
This paper reports \emph{measured} single-GPU inference performance for one model
(\textbf{Qwen3-8B-Q8\_0}) across three commodity edge accelerators --- AMD Ryzen AI Max+ 395
``Strix Halo'' (Radeon 8060S iGPU, gfx1151), NVIDIA GB10 ``Grace Blackwell'' (sm\_121), and
Apple M4 Max --- comparing \textbf{hanzo-engine} against \texttt{llama.cpp} on each, same model,
same metric (512-token prefill throughput \texttt{pp512}; single-stream decode throughput).
Every number here is measured on the named machine; AMD figures are carried from companion
paper~01 (measured 2026-06-08), NVIDIA and Apple figures were measured fresh on 2026-06-17.
We find two things, both with the numbers to back them. First, \textbf{hanzo-engine reaches
decode parity with \texttt{llama.cpp} on all three backends} (0.98$\times$/1.04$\times$/0.99$\times$),
while prefill trails by a consistent, backend-dependent margin (0.76--0.91$\times$). Second,
\textbf{the fastest edge box flips with the workload}: NVIDIA GB10 wins prefill (compute-bound),
Apple M4 Max wins decode by $\approx$2.5$\times$ (bandwidth-bound), and the GB10's Blackwell
decode path is the slowest of the three despite the strongest prefill. We are explicit about
what is \emph{not} here: a single model at a single quant, single-stream, no batching/serving
sweep, and AMD not re-run on this pass. Numbers that are projections or microbenchmarks are
labelled as such; there are none we invented.
\end{abstract}

\section{Why these three boxes}
The thesis of the edge-inference program is that the interesting target is not a datacenter
H100 but a \emph{commodity box a customer already owns}: a unified-memory APU, a small
Grace-Blackwell dev kit, or a MacBook. All three share LPDDR-class unified memory (no discrete
VRAM tier to hide behind), which makes decode --- a pure streaming read of every weight per
token --- brutally honest: decode throughput is memory bandwidth divided by model bytes, full
stop. Prefill, by contrast, is a compute-bound GEMM on the matrix/tensor cores. A backend can
win one regime and lose the other; this paper shows exactly that.

\paragraph{Hardware (as measured).}
\begin{itemize}[noitemsep]
  \item \textbf{AMD Strix Halo} --- Radeon 8060S iGPU, gfx1151 (RDNA3.5), 40 CU, native
        Windows + ROCm 7.1; measured streaming bandwidth $\approx$231 GB/s (paper~01).
  \item \textbf{NVIDIA GB10} --- Grace-Blackwell, sm\_121, CUDA 13, 121 GB unified; Linux aarch64.
  \item \textbf{Apple M4 Max} --- 137 GB unified, Metal; \texttt{llama.cpp} via the \texttt{BLAS,MTL} backend.
\end{itemize}

\section{Method}
One model throughout: \textbf{Qwen3-8B-Q8\_0} (36 layers, $n_\text{embd}$ 4096, $n_\text{ff}$ 12288,
32 query / 8 KV heads, head dim 128; $\approx$8.11 GiB resident). The 8-bit quant is the honest
``no accuracy excuse'' point --- it is the largest of the common quants, so decode is
bandwidth-bound and the comparison is purely about kernel/streaming efficiency, not about a
cheaper format hiding a slower engine.

\begin{itemize}[noitemsep]
  \item \textbf{Prefill} = \texttt{pp512}: tokens/s processing a 512-token prompt (compute-bound GEMM).
  \item \textbf{Decode} = single-stream token generation, tokens/s (bandwidth-bound mat-vec).
        For \texttt{llama.cpp} this is \texttt{tg128}; for hanzo-engine it is the bench command's
        128-token decode at depth 4. Both are single-stream token-generation rates at shallow
        KV-cache depth, treated as comparable decode throughput (not identical-length runs).
  \item All layers forced resident on the accelerator (\texttt{-ngl 99} for \texttt{llama.cpp};
        hanzo auto-maps all 36 layers to the GPU). \emph{This matters}: on a unified-memory box the
        device mapper can misread shared RAM and silently offload layers to the CPU, which tanks
        decode and invalidates the comparison.
  \item hanzo-engine built \texttt{--release} on each box (the NVIDIA binary was rebuilt release for
        this pass; a debug build is not a fair number). \texttt{llama.cpp} is the stock HIP/CUDA/Metal build.
\end{itemize}

\section{Results}
\begin{table}[h]
\centering
\caption{Qwen3-8B-Q8\_0, single GPU. Prefill = \texttt{pp512} t/s; Decode = single-stream t/s.
AMD from paper~01 (2026-06-08); NVIDIA + Apple measured 2026-06-17.}
\begin{tabular}{l rr c rr c}
\toprule
& \multicolumn{3}{c}{\textbf{Prefill (t/s)}} & \multicolumn{3}{c}{\textbf{Decode (t/s)}}\\
\cmidrule(lr){2-4}\cmidrule(lr){5-7}
Backend & llama.cpp & hanzo & ratio & llama.cpp & hanzo & ratio\\
\midrule
AMD Strix Halo (gfx1151) & 1176 & 1012 & 0.86$\times$ & 25.0 & 24.5 & 0.98$\times$\\
NVIDIA GB10 (sm\_121)     & 2155 & 1630 & 0.76$\times$ & 19.6 & 20.3 & \textbf{1.04$\times$}\\
Apple M4 Max (Metal)      &  873 &  792 & 0.91$\times$ & 50.4 & 49.9 & 0.99$\times$\\
\bottomrule
\end{tabular}
\end{table}

\noindent Two cross-cutting facts fall straight out of the table.

\paragraph{1. The fastest box depends on the workload.}
On \emph{prefill}, NVIDIA GB10 leads (2155 t/s, $1.8\times$ the APU, $2.5\times$ the M4 Max) ---
the Blackwell tensor cores dominate the compute-bound GEMM. On \emph{decode}, the order inverts
completely: \textbf{Apple M4 Max wins at 50.4 t/s}, about $2\times$ the Strix Halo and
$2.5\times$ the GB10. Decode is bandwidth-bound, and the M4 Max's wide unified memory
($\approx$50.4 t/s $\times$ 8.04 GB of weights $\approx$ 405 GB/s effective) simply moves weights faster.
Strikingly, the \textbf{GB10 --- the strongest prefill box --- is the \emph{slowest} at decode}
($\approx$19.6 t/s $\Rightarrow$ $\approx$160 GB/s effective): its LPDDR bandwidth and the
current \texttt{llama.cpp} Blackwell decode kernels leave it bandwidth-starved. A prompt-heavy /
RAG workload prefers the GB10; an interactive chat / long-generation workload prefers the M4 Max.
This is the central practical result for choosing edge hardware.

\paragraph{2. hanzo-engine is at decode parity with \texttt{llama.cpp} everywhere.}
hanzo's decode is 0.98$\times$, 1.04$\times$, and 0.99$\times$ of \texttt{llama.cpp} on AMD,
NVIDIA, and Apple respectively. We read all three as parity: these are best-of-3 single-stream
runs (one warmup) without a collected variance estimate, so the GB10's $1.04\times$ is within
run-to-run noise, not a measured win.
Since decode is the bandwidth wall, this says hanzo's per-format decode mat-vec is streaming at
the hardware ceiling on all three backends. Prefill trails by 0.76--0.91$\times$: smallest on
Metal (0.91$\times$) and largest on NVIDIA (0.76$\times$), where \texttt{llama.cpp}'s CUDA GEMM
is most heavily tuned. The prefill gap is real engineering work (matrix-core GEMM scheduling),
not a rounding error --- the same honest conclusion paper~01 reached on AMD.

\section{The edge/commodity angle: smaller models}
For latency-sensitive on-device use the 4B tier is the sweet spot. Measured \texttt{llama.cpp}
(2026-06-17), Qwen3-4B-Q8\_0:
\begin{table}[h]
\centering
\begin{tabular}{l rr}
\toprule
Backend & Prefill pp512 (t/s) & Decode (t/s)\\
\midrule
NVIDIA GB10 & 3913 & 36.0\\
Apple M4 Max & 1492 & 85.3\\
\bottomrule
\end{tabular}
\end{table}

\noindent The same pattern holds and sharpens: the M4 Max decodes the 4B at \textbf{85 t/s}
(comfortably real-time interactive), while the GB10 again leads prefill (3913 t/s) but trails
decode (36 t/s). Decode scales roughly inversely with model bytes (4B Q8 is $\approx$0.49$\times$
the 8B's bytes; M4 Max decode rises $50.4\to85.3$, $\approx$1.7$\times$), confirming the
bandwidth-bound model.

\section{Ported vs.\ novel; feasible vs.\ aspirational}
Consistent with paper~01, we separate what we copied from what is ours, and what is shipping from
what is roadmap.
\begin{itemize}[noitemsep]
  \item \textbf{Ported (deliberately):} the \texttt{mul\_mat\_q} integer-matrix-core prefill design
        and the GGUF block formats follow \texttt{llama.cpp}/\texttt{ggml}; the Metal path uses
        Apple's BLAS/MPS primitives as \texttt{llama.cpp} does.
  \item \textbf{Novel (hanzo):} the unified 1-bit-to-8-bit quant compute core (one kernel family
        serving the whole quant zoo, bit-exact vs the CPU reference), and a single multi-backend
        engine (CUDA/ROCm/Metal/Vulkan/CPU) behind one model loader, so the \emph{same} build runs
        every backend above. The cross-backend decode parity is evidence the core is not over-fit
        to one vendor.
  \item \textbf{FEASIBLE-now:} closing the NVIDIA prefill gap (0.76$\to$$\sim$0.9$\times$) via
        Blackwell-tuned GEMM tiles; a Blackwell decode kernel to lift the GB10 off its
        $\approx$160 GB/s effective floor.
  \item \textbf{ASPIRATIONAL:} beating \texttt{llama.cpp} prefill on any of these boxes; batched /
        multi-stream serving throughput (this paper is single-stream only).
\end{itemize}

\section{What this paper does \emph{not} claim}
Single model (Qwen3-8B-Q8\_0) at a single quant; single-stream decode (no concurrency/batching
sweep); \texttt{pp512}/128-token decode only (no long-context or KV-pressure regime); AMD numbers
carried from paper~01 rather than re-measured on this pass; and \texttt{llama.cpp} taken as the
stock build (not re-tuned per box). Accuracy is not evaluated here --- Q8\_0 is bit-faithful to
the reference, but no task-quality numbers are claimed. Everything reported is a measured
throughput on the named hardware; where a figure is derived (effective GB/s) it is computed from
the measured t/s and the model byte count, and labelled.

\section{Deployment guidance (for practitioners)}
The measurements translate into a concrete routing rule for anyone running models on hardware they
own:
\begin{table}[h]
\centering
\begin{tabular}{l l l}
\toprule
Workload & Best commodity box & Why\\
\midrule
RAG / long prompts / doc ingest & NVIDIA GB10 & prefill-bound; 2155 t/s, $1.8$--$2.5\times$ the others\\
Interactive chat / agents / long gen & Apple M4 Max & decode-bound; 50 t/s, $\approx2.5\times$ the GB10\\
Mixed / cost-first / "runs anything" & AMD Strix Halo APU & cheapest unified box; competitive both regimes\\
\bottomrule
\end{tabular}
\end{table}

\noindent Two operational points matter as much as the raw numbers. \textbf{(a) One binary, every
backend.} The same hanzo-engine build runs CUDA, ROCm, Metal, Vulkan, and CPU behind one model
loader --- there is no per-vendor inference stack to maintain, no rewrite to move a workload from a
Mac to a GB10 to an APU mini-PC. \textbf{(b) Commodity economics.} All three boxes are
consumer/prosumer hardware (a Mac, a small Grace-Blackwell dev kit, an APU mini-PC), not datacenter
GPUs, yet they run an 8B model at interactive decode and 800--2150 t/s prefill --- the edge is
genuinely viable, and the cheapest box (the APU) is within $\sim$2$\times$ of the others on both axes.

\section{Related work and contribution}
\texttt{llama.cpp} is the de-facto standard for edge/commodity LLM inference and is our baseline
throughout. Most published edge numbers, however, are either single-vendor (one GPU family) or
single-engine; cross-vendor comparisons typically vary both hardware \emph{and} software at once.
Our contribution is a controlled measurement: \emph{same} model, \emph{same} quant, \emph{same}
metric, the \emph{same} two engines, across three vendors --- which isolates hardware from software
and surfaces the workload-dependent hardware optimum (prefill$\to$NVIDIA, decode$\to$Apple) and the
fact that a single portable engine reaches decode parity with the hand-tuned baseline on all of them.

\section{Conclusion}
On commodity unified-memory accelerators, hanzo-engine delivers \texttt{llama.cpp}-class decode
(parity on AMD, NVIDIA, and Apple) from one multi-backend build, with a known and bounded prefill
gap. The more actionable result for anyone deploying at the edge is the workload split: buy/route
to \textbf{NVIDIA for prefill-heavy} pipelines and \textbf{Apple silicon for decode-heavy /
interactive} ones --- and do not assume the most expensive box wins decode, because here it loses.

\appendix
\section{Reproducibility}
\textbf{Hardware.} AMD Ryzen AI Max+ 395 ``Strix Halo'' (Radeon 8060S, gfx1151, 40 CU, $\approx$231 GB/s,
native Windows + ROCm 7.1); NVIDIA GB10 (Grace-Blackwell, sm\_121, CUDA 13, 121 GB unified, Linux
aarch64); Apple M4 Max (137 GB unified, Metal / \texttt{BLAS,MTL}).\\
\textbf{Software.} \texttt{llama.cpp} stock HIP/CUDA/Metal builds; hanzo-engine built \texttt{--release}
per box (git \texttt{8c50abb32}). Measurement passes: AMD 2026-06-08 (paper~01), NVIDIA + Apple 2026-06-17.\\
\textbf{Models.} \texttt{Qwen3-8B-Q8\_0} (8.11 GiB) and \texttt{Qwen3-4B-Q8\_0} (3.98 GiB), GGUF, from the
public Qwen3 releases.\\
\textbf{Commands.} Baseline: \texttt{llama-bench -m <model>.gguf -ngl 99 -p 512 -n 128}. hanzo:
\texttt{hanzo bench auto -m <dir> --format gguf -f <model>.gguf} (512-token prefill, 128-token decode,
1 warmup + 3 timed iterations).\\
\textbf{Protocol.} All layers resident on the accelerator (\texttt{-ngl 99} / hanzo auto-map of all 36
layers), single stream, quiet GPU. Effective bandwidth figures are derived as (decode t/s)\,$\times$\,(model
bytes) and labelled as derived, not measured directly.

\end{document}
